\chapter{Gauge fields in one dimension}\label{chap:gauge}

In one spatial dimension the electromagnetic field has no transverse polarizations: there are no photons. The
electric field is completely fixed by the charges through the Gauss law, and eliminating it leaves an
interaction between the charges. For bosonized fermions this interaction is quadratic, so the gauge theory stays
a system of coupled oscillators. This chapter derives the lattice Hamiltonian used by
\code{models.build_gauge_model} and its main physical consequences.

\section{Electric field and Gauss law}
In the temporal gauge $A_0 = 0$ the only field is $A_1(x)$ with conjugate momentum $-E(x)$, and the Hamiltonian
of the gauge field is $\frac12\int E^2$ (there is no magnetic field in one dimension). Physical states satisfy
the Gauss law
\begin{equation}\label{eq:gauss}
  \partial_x E = g\,\rho ,
\end{equation}
where $\rho$ is the charge density of the matter and $g$ the coupling. With the bosonized densities
\eqref{eq:dictionary} and charges \eqref{eq:va-charges},
\begin{equation}
  g\rho = \frac{g}{\sqrt\pi}\sum_i\bigl(Q^V_i\,\partial_x\varphi_i + Q^A_i\,\partial_x\vartheta_i\bigr)
  \quad\Rightarrow\quad
  E(x) = \frac{g}{\sqrt\pi}\sum_i\bigl(Q^V_i\varphi_i(x) + Q^A_i\vartheta_i(x)\bigr) + e(x) .
\end{equation}
The integration constant $e(x)$ is a c-number background: a constant $e = g\theta/2\pi$ (the $\theta$ angle) or
the field of static external charges (\cref{sec:probes}). The total Hamiltonian is
\begin{equation}\label{eq:H-gauge}
  H = \int dx\,\Bigl[\tfrac12\sum_i\bigl(\Pi_i^2 + (\partial_x\varphi_i)^2\bigr) + \tfrac12 E^2\Bigr].
\end{equation}
With several gauge fields $U(1)^k$ there is one $E_a$ with coupling $g_a$ and charges $Q^V_a, Q^A_a$ per gauge
field, and the electric energy is $\frac12\sum_a\int E_a^2$.

\paragraph{On the lattice.} With $\vartheta_{i,j} = \sum_{k\leq j}P_{i,k}$ (the discrete version of
$\int^x\Pi$) the electric field on site $j$ is a linear function of $\X$,
\begin{equation}
  E_j = (C\X)_j + e_j,\qquad
  E_j = \frac{g}{\sqrt\pi}\sum_i\Bigl(Q^V_i\varphi_{i,j} + Q^A_i\sum_{k\leq j}P_{i,k}\Bigr) + e_j ,
\end{equation}
and the electric energy $\frac a2\sum_jE_j^2$ adds $aC^TC$ to the matrix $M$ (plus a linear term and a constant
if $e\neq0$). On the ring the zero modes enter as well:
$\varphi_{i,j} = \tilde\varphi_{i,j} + W_ij/N$ and $\vartheta_{i,j} = \bar\vartheta_i + \sum_{k\leq j}P_{i,k}$.

\section{The Schwinger model}
The Schwinger model~\cite{Schwinger1962,LowensteinSwieca1971} is a single Dirac fermion with charge $1$
($Q_L = Q_R = 1$, so $Q^V = 1$, $Q^A = 0$). Then $E = (g/\sqrt\pi)\varphi$ and \eqref{eq:H-gauge} becomes
\begin{equation}
  H = \int dx\,\Bigl[\tfrac12\Pi^2 + \tfrac12(\partial_x\varphi)^2 + \frac{g^2}{2\pi}\varphi^2\Bigr]:
\end{equation}
a free boson of mass
\begin{equation}\label{eq:schwinger-mass}
  \mu = \frac{g}{\sqrt\pi} .
\end{equation}
The massless fermions have disappeared from the spectrum; the only particle is a massive neutral boson, the
``meson''. Its mass is generated non-perturbatively by the gauge field.

\paragraph{General charges.} For one gauge field with charges $Q_L,Q_R$ (anomaly free, \cref{chap:anomalies})
exactly one combination of the bosons becomes massive, with
\begin{equation}\label{eq:meson-mass}
  \mu^2 = \frac{g^2}{\pi}\sum_i\bigl[(Q^V_i)^2 + (Q^A_i)^2\bigr] = \frac{g^2}{2\pi}\sum\bigl(Q_L^2 + Q_R^2\bigr),
\end{equation}
and the remaining $F-1$ combinations stay massless. For $U(1)^k$ the masses are the square roots of the
eigenvalues of $\mu^2_{ab} = (g_ag_b/2\pi)(Q_{L,a}\cdot Q_{L,b} + Q_{R,a}\cdot Q_{R,b})$
(\code{ChiralTheory.meson_masses}).

\section{Spectra on the open chain}
On the open chain the predicted spectrum is one branch $\sqrt{\mu^2 + k_n^2}$ and $F-1$ massless branches $k_n$,
with the Dirichlet momenta \eqref{eq:dirichlet}. For vector charges the electric field acts locally on $\varphi$
and the prediction is exact on the lattice. Axial charges act on $\vartheta$, a sum over all sites to the left,
whose boundary conditions differ from Dirichlet: the momenta become $k_n = (n+\delta)\pi/L$ with a phase shift
$\delta$, and the prediction holds only in the bulk. \Cref{fig:spectra_open_chain,fig:finite_size_scaling} show
both effects; the deviation for axial charges falls like $1/N$.

\latsimfigure{spectra_open_chain}
\latsimfigure{finite_size_scaling}

\paragraph{Continuum limit.} At fixed physical length $L$ the lowest frequency of the Schwinger model approaches
$\sqrt{\mu^2 + (\pi/L)^2}$ as $a\to0$, with the $(ka)^2$ lattice correction of \cref{chap:lattice}
(\cref{fig:continuum_limit}).

\latsimfigure{continuum_limit}

\section{Measuring the mass: the electric-field correlator}
The electric field is proportional to the massive boson, so its correlations measure $\mu$. For a free boson of
mass $\mu$ the equal-time correlator in the ground state is
\begin{equation}
  \expval{\Phi(x)\Phi(0)} = \int\frac{dk}{2\pi}\,\frac{e^{ikx}}{2\sqrt{k^2+\mu^2}} = \frac{1}{2\pi}K_0(\mu|x|),
\end{equation}
with the modified Bessel function $K_0$, which decays like $e^{-\mu|x|}$. The library computes
$\expval{E_jE_k} = (C\sigma C^T)_{jk}$ from the covariance matrix \eqref{eq:sigma} and fits $A K_0(\mu r) + B$
in a window away from the boundaries (\code{analysis.correlators.meson_mass_from_E} with an explicit
\code{FitWindow}). The offset $B$ is a uniform fluctuation of the background field caused by the open boundary,
non-zero for axial charges. The fitted masses agree with \eqref{eq:meson-mass} to about $10^{-3}$
(\cref{fig:correlator_fits}).

\latsimfigure{correlator_fits}

\section{Probe charges: screening and confinement}\label{sec:probes}
Place static external charges $+q$ at $x=0$ and $-q$ at $x = d$. By the Gauss law they create the background
$e(x) = gq$ for $0<x<d$. The ground-state energy shift is the potential $V(d)$ between them. For the Schwinger
model the Hamiltonian is $\int[\frac12\Pi^2 + \frac12(\partial\varphi)^2 + \frac12(\mu\varphi + e)^2]$, and
minimizing over $\varphi$ in Fourier space gives
\begin{equation}
  V(d) = \frac12\int\frac{dk}{2\pi}\,|e_k|^2\frac{k^2}{k^2+\mu^2}
       = g^2q^2\int\frac{dk}{2\pi}\frac{1 - \cos kd}{k^2+\mu^2}
       = \frac{g^2q^2}{2\mu}\bigl(1 - e^{-\mu d}\bigr).
\end{equation}
The potential saturates: the massless fermions \emph{screen} any charge, even a fractional one~\cite{CJS1975}.
With $\mu = g/\sqrt\pi$ the prefactor is $g^2q^2/2\mu = \frac{\sqrt\pi}{2}gq^2$, which is exactly the flat-space
result quoted in~\cite{BharadwajIsenKang2026}, eq.~(3.88); the lattice reproduces it to $5\times10^{-4}$.
Without charged matter ($\mu = 0$) the same formula gives $V(d) = g^2q^2d/2$: a constant electric field between
the charges, linear \emph{confinement} with string tension $g^2q^2/2$. On the lattice the potential is the
background energy $\min_{\X}[\frac12\X^TM\X + \vect{b}^T\X] + \text{const} = \text{const} - \frac12\vect{b}^TM^+\vect{b}$
(\code{api.screening_potential}); \cref{fig:probe_potentials} shows both regimes.

\latsimfigure{probe_potentials}

\section{The theta angle}
A constant background field $e = g\theta/2\pi$ is the $\theta$ angle of the theory. With massless fermions it is
unphysical: in the Schwinger model it is removed by the shift $\varphi\to\varphi - \sqrt\pi\,\theta/2\pi$, which
corresponds to a chiral rotation of the fermions. On the ring this shift is a zero mode and costs nothing, so the
background energy vanishes (numerically $10^{-15}$). On the open chain the Dirichlet boundary fixes $\varphi$ at the
ends and a small boundary energy remains. A physical $\theta$ dependence requires massive fermions, which are
beyond the Gaussian methods of this release.

\section{In code}
\begin{python}
from latsimpkg import schwinger, ChiralTheory, spectrum, screening_potential

sp = spectrum(schwinger(g=0.04), N=300)          # open chain
sp.max_deviation                                 # ~1e-14: exact for vector charges
pot = screening_potential(schwinger(g=0.1), N=400, distances=range(2, 120, 6), q=0.5)
pot.meta["regime"], pot.max_rel_deviation        # 'screening', ~5e-4
\end{python}
